\section{The entire original DVR intersection in the prescribed fraction field}
\label{sec:regular-dvr-intersection}

Let $R$ be the original regular domain and let $F$ be its original
prescribed fraction field. For each original prime $P$ whose actual
localization $R_P$ is a discrete valuation ring, view that localization
as its actual original subalgebra in $F$.

\begin{theorem}[Every original coefficient in all original DVR localizations is regular]
\label{sc0-regular-dvr-intersection:original-ring}
\lean{EnvelopeProperty.exists_originalRing_of_mem_all_regularDVR_localizations}
\leanok
If an original $q\in F$ belongs to EVERY such original localization
subalgebra, there exists an original $a\in R$ whose original image
in the SAME prescribed $F$ equals $q$. No height-one coverage,
normality, intersection equality or coefficient-comparison premise
is supplied.
\end{theorem}
\begin{proof}
Write the SAME original coefficient as $a/b$ in $F$ with original
$b\ne0$. The actual principal-ideal detection theorem for regular
rings reduces $a\in(b)$ to detection at all original minimal primes
of $(b)$. Their actual localizations are discrete valuation rings,
as already derived from regularity and the genuine principal-prime
argument, so the hypothesis applies to each of these SAME primes.
At such a prime write $q=c/s$ with the original $s$ outside that
prime. Equality in the SAME prescribed fraction field gives the
original cross-multiplication $sa=bc$, by injectivity of the
original fraction-field map. This is exactly the required localized
membership in the original principal ideal. Principal-ideal
detection therefore gives $a=bc$ for some original $c\in R$.
Cancellation in the SAME original $F$ yields $q=c$, as required.
\end{proof}

\begin{theorem}[Whole original DVR intersection]
\label{sc0-regular-dvr-intersection:whole-subalgebra}
\lean{EnvelopeProperty.iInf_regularDVR_localizations_eq_bot_in_fractionField}
\leanok
The ENTIRE intersection of those actual original DVR localization
subalgebras in the SAME prescribed $F$ is precisely the original
$R$-subalgebra, the bottom subalgebra of $F$ over $R$.
\end{theorem}
\begin{proof}
An element of the entire original intersection belongs to every
original localization. The preceding theorem expresses it as the
image of an actual original ring element, hence as an element of
the bottom subalgebra. Conversely the bottom subalgebra is contained
in every original $R$-subalgebra, proving the full object equality.
\uses{sc0-regular-dvr-intersection:original-ring}
\end{proof}

This preserves the original specified field and ALL its original
coefficient maps. It is the algebraic regularity criterion for the
actual proper valuative-lift canonical-section argument. Deriving
that every original rounded canonical coefficient belongs to these
localizations, and then the entire canonical-module pushforward,
remain subsequent obligations. No vanishing claim follows here.
